\section{场景组合细节}
\label{sec:supp-composition}

\paragraph{状态驱动渲染。}

在线步 $k$ 的状态 $s_k$ 指定自车位姿、活动资产、刚体变换，以及人体根位姿与关节姿态。刚体资产的高斯中心和协方差变换为

\begin{equation}
\boldsymbol\mu^w=\mathbf R\boldsymbol\mu^a+\mathbf t,
\qquad
\boldsymbol\Sigma^w=\mathbf R\boldsymbol\Sigma^a\mathbf R^\top.
\label{eq:rigid}
\end{equation}

人体使用加权蒙皮变换 $\mathbf A_i=\sum_j w_{ij}\mathbf P_j\mathbf B_j^{-1}$，其中 $\mathbf B_j$ 和 $\mathbf P_j$ 是根坐标系中的静止与当前全局关节矩阵。变换后的中心为 $\bar{\boldsymbol\mu}_i^w=\mathbf T^w_h\mathbf A_i\bar{\boldsymbol\mu}_i^h$。组合局部线性映射 $\mathbf J_i$ 将协方差变换为 $\mathbf J_i\boldsymbol\Sigma_i\mathbf J_i^\top$，其逆 $\mathbf J_i^{-1}$ 将观察方向映回规范坐标，并在外观求值前归一化。上述逆映射使用非奇异蒙皮变换。根位姿控制整体放置，根关节保持静止变换。

渲染器合成背景、活动资产和有效天空。外部策略与动力学可推进状态：

\begin{align}
o_k&=\mathcal R(\operatorname{Compose}(\mathcal W,s_k),\mathcal C_k),\\
s_{k+1}&=F_{\rm dyn}(s_k,\pi(o_k),\Delta t).
\label{eq:simulation}
\end{align}

相机与参与者状态变化时复用同一背景和规范资产。道路网格、碰撞代理与视觉组件采用相同坐标，可由物理仿真器驱动其运动。
